
%
%
%
%
%
%
\chapter{Natural Semantics and Structural Operational Semantics: Formalization and Equivalence}
\label{formal_ns_sos_equiv}
First steps before reasoning wheter or not Natural Semantics and Structural Operational Semantics are equivilent is to create them.
and to create them we also need a langugae to apply them to, in our case this will be \While.


%
%
%
\section{Formalizing \While}
\label{formal_ns_sos_equiv:sec:while:formalization}

Due to Lean 4 being a programming language,
it is needed to define the terms before using them,
this will lead to some textbook definitions to come out of order.

We will begin by mechanizing the syntactic categories defined in the textbook in chapter 1.2.
The left-hand-side contains the Lean 4 code which represents the textbook definition given to the left.

The syntactic categories for \textit{While} are defined using BNF notation. We map numerals ($n$), variables ($x$), and expressions ($a, b$) to inductive sets.

\begin{paracol}{2}
    
    % --- SECTION: SYNTACTIC CATEGORIES ---
    \begin{leftcolumn}
        \centering \textbf{Mechanized (Lean 4)}
    \end{leftcolumn}
    \begin{rightcolumn}
        \centering \textbf{Pen-and-Paper (Textbook)} \\\noindent
    \end{rightcolumn}

    \switchcolumn[1]*[]
    \begin{leftcolumn}
        \begin{Lean_Code}
inductive Num where
  | zero : Num
  | one  : Num
  | succ0 (n : Num) : Num
  | succ1 (n : Num) : Num
        \end{Lean_Code}
    \end{leftcolumn}
    \begin{rightcolumn}
        \(\displaystyle
        \begin{array}{lcl}
            n &::=& 0 \\ 
              &\mid& 1 \\ 
              &\mid& n\ 1 \\ 
              &\mid& n\ 0
        \end{array}
        \)
    \end{rightcolumn}
\end{paracol}
\vspace{1em}
As shown above, the translation of inductive categories are easy to translate into Lean 4.
Using the keyword inductive starts the definition, where the then give it a name, in this case \emph{Num}, we then continue with the keyword \emph{where}
which marks the start of the case by case definition.
Next we give our inductive case a name, then follow it by its type, in this case \emph{Num}, ... fix this
On the last two lines we also include a variable n, this marks a recursive relationship to itself.

We continue with arithmetic expressions:
\begin{paracol}{2}
    \begin{leftcolumn}
        \centering \textbf{Mechanized (Lean 4)}
    \end{leftcolumn}
    \begin{rightcolumn}
        \centering \textbf{Pen-and-Paper (Textbook)} \\\noindent
    \end{rightcolumn}

    \switchcolumn[1]*[]
    \begin{leftcolumn}
        \begin{Lean_Code}
inductive Aexp where
  | num (n : Num)
  | var (x : Var)
  | add (a1 a2 : Aexp)
  | mul (a1 a2 : Aexp)
  | sub (a1 a2 : Aexp)
        \end{Lean_Code}
    \end{leftcolumn}
    \begin{rightcolumn}
        \(\displaystyle
        \begin{array}{lcl}
            a &::=& n \\ 
              &\mid& x \\ 
              &\mid& a_1 + a_2 \\ 
              &\mid& a_1 \star a_2 \\ 
              &\mid& a_1 - a_2 \\
        \end{array}
        \)
    \end{rightcolumn}
\end{paracol}
\vspace{1em}
Here we can see a more clear usage of other categories used to define this new category.
In this case we use \emph{Num}, \emph{Var} and \emph{Aexp}. notice also that we are able to give the same type to multiple variables at once as seen in the binary operations.

We continue with boolean expressions:
\begin{paracol}{2}
    \begin{leftcolumn}
        \begin{Lean_Code}
inductive Bexp where
  | true
  | false
  | eq  (a1 a2 : Aexp)
  | le  (a1 a2 : Aexp)
  | not (b1 : Bexp)
  | and (b1 b2 : Bexp)
        \end{Lean_Code}
    \end{leftcolumn}
    \begin{rightcolumn}
        \(\displaystyle
        \begin{array}{lcl}
            b &::=& \mathbf{true} \\ 
              &\mid& \mathbf{false} \\ 
              &\mid& a_1 = a_2 \\ 
              &\mid& a_1 \leq a_2 \\ 
              &\mid& \neg b_1 \\ 
              &\mid& b_1 \land b_2
        \end{array}
        \)
    \end{rightcolumn}
\end{paracol}
\vspace{1em}
This is more of the same, but one intresting decision is to use the mathmatical true, and false instead of the \textbf{tt} and \textbf{tt} and \textbf{ff} used in the literature.
This decision was made to be able to use the builtin methods in Lean to be able to reason about boolean algebra, such as negation of a negation etc.

Lastly we define our statments:
\begin{paracol}{2}
    \begin{leftcolumn}
        \begin{Lean_Code}
inductive Stmt where
  | ass  (x : Var)  (a : Aexp)
  | skip
  | composition     (S@$_{1}$@ S@$_{2}$@ : Stmt)
  | cond (b : Bexp) (S@$_{1}$@ S@$_{2}$@ : Stmt)
  | loop (b : Bexp) (S@$_{1}$@ : Stmt)
        \end{Lean_Code}
    \end{leftcolumn}
    \begin{rightcolumn}
        \(\displaystyle
        \begin{array}{lcl}
    S &::=& x := a \\
      &\mid& \mathtt{skip} \\
      &\mid& S_1; S_2 \\
      &\mid& \mathtt{if}\ b\ \mathtt{then}\ S_1\ \mathtt{else}\ S_2 \\
      &\mid& \mathtt{while}\ b\ \mathtt{then}\ S
        \end{array}
        \)
    \end{rightcolumn}
\end{paracol}
\vspace{1em}
Here we were meet with some namespace collisions as the keywords \emph{if} and \emph{while} are used already, leading us to change it to
\emph{cond} to represent a condition, in this case that would be the \emph{if}-statement, we also use \emph{loop} to represent looping statements, in this case the \emph{while}-statement.

The \textbf{State} represents memory as a total function from variables to integers.
A definition for substitution is also needed.
\begin{paracol}{2}
    % --- SECTION: SEMANTIC DOMAINS ---
    \begin{leftcolumn}
        \centering \textbf{Mechanized (Lean 4)}
    \end{leftcolumn}
    \begin{rightcolumn}
        \centering \textbf{Pen-and-Paper (Textbook)} \\\noindent
    \end{rightcolumn}

    \switchcolumn[1]*[]
    \begin{leftcolumn}
        \begin{Lean_Code}
def State := Var @\to@ @ \(\mathbb{Z}\)@
        \end{Lean_Code}
    \end{leftcolumn}
    \begin{rightcolumn}
        \(\displaystyle s : \mathbf{Var} \to \mathbb{Z} \)
    \end{rightcolumn}

    \switchcolumn[1]*[]
    \begin{leftcolumn}
        \begin{Lean_Code}
def assign (s : State) (x : Var) (z : @\(\mathbb{Z}\)@) : State :=
  fun (v: Var) => 
    if v = x then z else s v

notation s "[" x " @\color{purple}\mapsto@ " z "]" @\leanbreak@=> assign s x z
        \end{Lean_Code}
    \end{leftcolumn}
    \begin{rightcolumn}
        \(\displaystyle
        s[x \mapsto v]y = \begin{cases} 
        v & \text{if } y = x \\ 
        s\,y & \text{if } y \neq x 
        \end{cases}
        \)
    \end{rightcolumn}

\end{paracol}
\vspace{1em}
For a more detailed breakdown see \autoref{app:lean:state} in the appendix.
The character "\(\mbox{\textcolor{gray}{$\hookrightarrow$}}\)" denotes a line break due to length limitations in the report, in the code this line is unbroken.

Notice that the notation is after the definition, this is due to a technicality that the function \emph{assign} is not defined untill the end of its definition,
which means that anything that is referring to the function \emph{assign} must be placed underneath its definition to be valid.
Now that we have defined the syntactic categories we can provide the semantic functions and equations.
We skip defining the syntactic domains as they are implicitly defined in the functions to come.

Each function is first  defined as a relation from one domain to another, and then the implementation of the function is given using the different semantic clauses.
We will begin by providing the definition for Numerals (\Num):
\begin{paracol}{2}
    % --- SECTION: SEMANTIC FUNCTIONS ---
    \begin{leftcolumn}
        \centering \textbf{Mechanized (Lean 4)}
    \end{leftcolumn}
    \begin{rightcolumn}
        \centering \textbf{Pen-and-Paper (Textbook)} \\\noindent
    \end{rightcolumn}

    \switchcolumn[1]*[]
    \begin{leftcolumn}
        \begin{Lean_Code}
def Num_to_Z : Num @\to@ @\(\mathbb{Z}\)@
        \end{Lean_Code}
    \end{leftcolumn}
    \begin{rightcolumn}
        \[
        \mathcal{N} : Num \to \mathbb{Z}
        \]
    \end{rightcolumn}
\end{paracol}
\vspace{1em}
The start of defining the function using our categories is the same,
the only difference here is the usage of keyword \emph{def}.

To define how the function works we in the lean case follow by doing a case by case break down in the same definitions,
compared to the textbook which can defer the definition until later.
\begin{paracol}{2}
    \begin{leftcolumn}
        \centering \textbf{Mechanized (Lean 4)}
    \end{leftcolumn}
    \begin{rightcolumn}
        \centering \textbf{Pen-and-Paper (Textbook)} \\\noindent
    \end{rightcolumn}

    \switchcolumn[1]*[]
    \begin{leftcolumn}
        \begin{Lean_Code}
def Num_to_Z : Num @\to@ @\(\mathbb{Z}\)@
  | Num.zero    => 0
  | Num.one     => 1
  | Num.succ0 n => 2 * (Num_to_Z n)
  | Num.succ1 n => 2 * (Num_to_Z n) + 1

notation "@\color{purple}\(\mathcal{N}[\![\)@" n "@\color{purple}\(]\!]\)@" => Num_to_Z n
        \end{Lean_Code}
    \end{leftcolumn}
    \begin{rightcolumn}
        \[
        \begin{array}{rcl}
            \mathcal{N}[\![0]\!]s &=& 0 \\[0.6em]
            \mathcal{N}[\![1]\!]s &=& 1 \\[0.6em]
            \mathcal{N}[\![n\ 0]\!]s &=& 2 \cdot \mathcal{N}[\![n]\!] \\[0.6em]
            \mathcal{N}[\![n\ 1]\!]s &=& 2 \cdot \mathcal{N}[\![n]\!] + 1
        \end{array}
        \]
    \end{rightcolumn}
\end{paracol}
\vspace{1em}
Here we create a function, denoted by the \emph{def} keyword, in Lean this defaults to being a Total Function.
By defining how the function acts for all inductive cases of the type we are working with (in this case \emph{Num}) we ensure that it is total.
The case by case implementation is trivial and directly lifted from the literature.

Next we will provide the definition for arithmetic expressions (\Aexp):
\begin{paracol}{2}
    \begin{leftcolumn}
        \centering \textbf{Mechanized (Lean 4)}
    \end{leftcolumn}
    \begin{rightcolumn}
        \centering \textbf{Pen-and-Paper (Textbook)} \\\noindent
    \end{rightcolumn}

    \switchcolumn[1]*[]
    \begin{leftcolumn}
        \begin{Lean_Code}
def Aexp_eval : Aexp @\to@ (State @\to@ @\(\mathbb{Z}\)@)
  | Aexp.num n, _ => @ \(\mathcal{N}[\![n]\!]\)@
  | Aexp.var x, s => s x
  | Aexp.add @\(a_1\)@ @\(a_2\)@, s => Aexp_eval @\(a_1\)@ s + Aexp_eval @\(a_2\)@ s
  | Aexp.mul @\(a_1\)@ @\(a_2\)@, s => Aexp_eval @\(a_1\)@ s * Aexp_eval @\(a_2\)@ s
  | Aexp.sub @\(a_1\)@ @\(a_2\)@, s => Aexp_eval @\(a_1\)@ s - Aexp_eval @\(a_2\)@ s

notation "@\color{purple}\(\mathcal{A}[\![\)@" a "@\color{purple}\(]\!]\)@" @\leanbreak@=> Aexp_eval a
        \end{Lean_Code}
    \end{leftcolumn}
    \begin{rightcolumn}
        \[
        \mathcal{A} : Aexp \to (State \to \mathbb{Z})
        \]
        and is defined by
        \[
        \begin{array}{rcl}
            \mathcal{A}[\![n]\!]s &=& \mathcal{N}[\![n]\!] \\[0.6em]
            \mathcal{A}[\![x]\!]s &=& s \, x \\[0.6em]
            \mathcal{A}[\![a_1 + a_2]\!]s &=& \mathcal{A}[\![a_1]\!]s + \mathcal{A}[\![a_2]\!]s \\[0.6em]
            \mathcal{A}[\![a_1 \star a_2]\!]s &=& \mathcal{A}[\![a_1]\!]s \cdot \mathcal{A}[\![a_2]\!]s \\[0.6em]
            \mathcal{A}[\![a_1 - a_2]\!]s &=& \mathcal{A}[\![a_1]\!]s - \mathcal{A}[\![a_2]\!]s
        \end{array}
        \]
    \end{rightcolumn}
\end{paracol}
\vspace{1em}
In the above snippet we have the first underscore "\_" used, this is a reserved character used to let Lean 4 itself infer
what variable to use, in this case it is used instead of writing out the state \emph{s}, which the rule does not care for.
Lastly we will provide the definition for boolean expressions (\Bexp):
\begin{paracol}{2}
    \begin{leftcolumn}
        \centering \textbf{Mechanized (Lean 4)}
    \end{leftcolumn}
    \begin{rightcolumn}
        \centering \textbf{Pen-and-Paper (Textbook)} \\\noindent
    \end{rightcolumn}

    \switchcolumn[1]*[]
    \begin{leftcolumn}
        \begin{Lean_Code}
def Bexp_eval : Bexp @\to@ (State @\to@ @\(\mathbb{Z}\)@)
  | Bexp.@true@, _ => true
  | Bexp.@false@, _ => false
  | Bexp.eq @\(a_1\)@ @\(a_2\)@, s => Aexp_eval @\(a_1\)@ s == Aexp_eval @\(a_2\)@ s
  | Bexp.le @\(a_1\)@ @\(a_2\)@, s => Aexp_eval @\(a_1\)@ s @\leq@ Aexp_eval @\(a_2\)@ s
  | Bexp.not @\(b\)@, s => @\leanbreak@@\neg@ Bexp_eval @\(b\)@ s
  | Bexp.and @\(b_1\)@ @\(b_2\)@, s => Bexp_eval @\(b_1\)@ s @\land@ Bexp_eval @\(b_2\)@ s

notation "@\color{purple}\(\mathcal{B}[\![\)@" b "@\color{purple}\(]\!]\)@" @\leanbreak@=> Bexp_eval b
        \end{Lean_Code}
    \end{leftcolumn}
    \begin{rightcolumn}
        \[
        \mathcal{B} : Bexp \to (State \to \mathbb{Z})
        \]
        \[
        \begin{array}{rcl}
            \mathcal{B}[\![true]\!]s &=& \textbf{tt} \\[0.6em]
            \mathcal{B}[\![false]\!]s &=& \textbf{ff} \\[0.6em]
            \mathcal{B}[\![a_1 = a_2]\!]s &=& \begin{cases} 
        \textbf{tt} & \text{if } \mathcal{A}[\![a_1]\!]s = \mathcal{A}[\![a_2]\!]s \\ 
        \textbf{ff} & \text{if } \mathcal{A}[\![a_1]\!]s \neq \mathcal{A}[\![a_2]\!]s
        \end{cases} \\[1.2em]
            \mathcal{B}[\![a_1 \leq a_2]\!]s &=& \begin{cases} 
        \textbf{tt} & \text{if } \mathcal{A}[\![a_1]\!]s \leq \mathcal{A}[\![a_2]\!]s  \\ 
        \textbf{ff} & \text{if } \mathcal{A}[\![a_1]\!]s > \mathcal{A}[\![a_2]\!]s 
        \end{cases} \\[1.2em]
            \mathcal{B}[\![\neg b]\!]s &=& \begin{cases} 
        \textbf{tt} & \text{if } \neg \mathcal{B}[\![b]\!]s \\ 
        \textbf{ff} & \text{if } \mathcal{B}[\![b]\!]s
        \end{cases} \\[1.2em]
            \mathcal{B}[\![b_1 \land b_2]\!]s &=& \begin{cases} 
        \textbf{tt} & \text{if } \mathcal{B}[\![b_1]\!]s = \textbf{tt} \text{ and } \mathcal{B}[\![b_2]\!]s = \textbf{tt} \\ 
        \textbf{ff} & \text{if } \mathcal{B}[\![b_1]\!]s = \textbf{ff} \text{ or } \mathcal{B}[\![b_2]\!]s = \textbf{ff} \\ 
        \end{cases}
        \end{array}
        \]
    \end{rightcolumn}
\end{paracol}
\vspace{1em}
In addition to the above definitions we will also need some extra proofs to allow us to reason about states.
First we need a proof to prove that
\[
    s[x \mapsto z_1][x \mapsto z_2] = s[x \mapsto z_2]
\]
This proof would allow us to remove redundant assignments which will be useful when reasoning about repeated assignments in things like loops.
There was no pen-and-paper proof provided for this property in the textbook.
But one could imagine the proof to look something  like this:
\todo{Fill in a pen-and-paper proof here please}

\noindent
In Lean this proof is constructed  using the following:
\begin{Lean_Code}
theorem assign_shadow (s : State) (x : Var) (z1 z2 : @\(\mathbb{Z}\)@) :
  s[x @\mapsto@ z1][x @\mapsto@ z2] = s[x @\mapsto@ z2] := by @\dots@
\end{Lean_Code}

The next property is commutative property of distinct substitutions.
A proof of this property needs to effectively prove the statement:
\[
    s[x \mapsto z_1][y \mapsto z_2] = s[y \mapsto z_2][x \mapsto z_1], y \neq x
\]
This proof would allow us to remove even more redundant assignments which will be useful when reasoning about repeated assignments in things like loops.
There was no pen-and-paper proof provided for this property in the textbook.
But one could imagine the proof to look something  like this:
\todo{Fill in a pen-and-paper proof here please}


\noindent
In Lean this proof is constructed  using the following theorem:
\begin{Lean_Code}
theorem assign_comm (s : State) (x y : Var) @\leanbreak@ (z1 z2 : @\(\mathbb{Z}\)@) (h : x @\neq@ y) :
  s[x @\mapsto@ z1][y @\mapsto@ z2] = s[y @\mapsto@ z2][x @\mapsto@ z1] := by @\dots@
\end{Lean_Code}
Full proofs for the two auxiliary theorems can be found in \autoref{app:lean:state} in the appendix.
By adding  these proofs it becomes trivial to prove a statement such as:
\[
    s[x \mapsto 3][x \mapsto 2][x \mapsto 1][x \mapsto 0] = s[x \mapsto 0]
\]
and things like
\[
    s[x \mapsto 3][y \mapsto 1][x \mapsto 2][y \mapsto 3][x \mapsto 1][y \mapsto 6][x \mapsto 0][y \mapsto 6] = s[x \mapsto 0][y \mapsto 6]
\]
which otherwise might not be obvius for the reader, and definitly not for the Lean 4 complier.

%
%
%
\section{Formalizing Natural Semantics (Big-Step)}
\label{formal_ns_sos_equiv:sec:ns:formalization}
Now that we have a mechanized definition of the \While language we can proceed by defining the actual semantics of running a \While language program. 
We will start by defining the natural operational semantics often referd to as "big step semantics".
When it comes to operational semantics we use deduction tree syntax to provide rules and axionoms.

\begin{paracol}{2}
    \begin{leftcolumn}
        \centering \textbf{Mechanized (Lean 4)}
    \end{leftcolumn}
    \begin{rightcolumn}
        \centering \textbf{Pen-and-Paper (Textbook)} \\\noindent
    \end{rightcolumn}

    \switchcolumn[1]*[]
    \begin{leftcolumn}
        \begin{Lean_Code}
inductive big_step : Stmt @\to@ State @\to@ State @\to@ Prop @\leanbreak@ where
        \end{Lean_Code}
    \end{leftcolumn}
    \begin{rightcolumn}
    \end{rightcolumn}
    \vspace{2em}
    \switchcolumn[1]*[]
    \begin{leftcolumn}
        \begin{Lean_Code}
  | ass {s x a} :
      big_step (Stmt.ass x a) s (s[x @\mapsto@ @\(\mathcal{A}[\![a]\!]\)@ s])
\end{Lean_Code}
    \end{leftcolumn}
    \begin{rightcolumn}
        \vspace{1.1em}
            \(
\begin{prooftree}
  \hypo{} % Empty hypothesis for an axiom
  \infer[left label=[ass$_{\text{ns}}$]]1{ \langle x := a, s \rangle \to s[x \mapsto \mathcal{A}[\![a]\!] s]}
\end{prooftree}
            \)
    \end{rightcolumn}
    \vspace{1em}
    \switchcolumn[1]*[]
    \begin{leftcolumn}
        \begin{Lean_Code}
  | skip {s} :
      big_step Stmt.skip s s
        \end{Lean_Code}
    \end{leftcolumn}
    \begin{rightcolumn}
        \vspace{1.1em}
            \(
\begin{prooftree}
  \hypo{} % Empty hypothesis for an axiom
  \infer[left label=[skip$_{\text{ns}}$]]1{ \langle \mathtt{skip}, s \rangle \to s}
\end{prooftree}
            \)
    \end{rightcolumn}
    \vspace{1em}
    \switchcolumn[1]*[]
    \begin{leftcolumn}
        \begin{Lean_Code}
  | comp {S@$_{1}$@ S@$_{2}$@ s s' s''}
    (h_left  : big_step S@$_{1}$@ s  s')
    (h_right : big_step S@$_{2}$@ s' s'') :
      big_step (Stmt.composition S@$_{1}$@ S@$_{2}$@) s s''
        \end{Lean_Code}
    \end{leftcolumn}
    \begin{rightcolumn}
        \vspace{1.1em}
            \(
\begin{prooftree}
  \hypo{\langle S_1, s \rangle \to s'}
  \hypo{\langle S_2, s' \rangle \to s''}
  \infer[left label=[comp$_{\text{ns}}$]]2{ \langle S_1;\ S_2, s \rangle \to s''}
\end{prooftree}
            \)
    \end{rightcolumn}
    \vspace{3em}
    \switchcolumn[1]*[]
    \begin{leftcolumn}
        \begin{Lean_Code}
  | if_true {b S@$_{1}$@ S@$_{2}$@ s s'} :
    (h_cond : @\(\mathcal{B}[\![b]\!]\)@ s = true) @\to@
    (h_then : big_step S@$_{1}$@ s s') @\to@
      big_step (Stmt.cond b S@$_{1}$@ S@$_{2}$@) s s'
        \end{Lean_Code}
    \end{leftcolumn}
    \begin{rightcolumn}
        \vspace{1.1em}
            \(
\begin{prooftree}
  \hypo{\langle S_1, s \rangle \to s'}
  \infer[
    left label={[if$_{\text{ns}}^{\text{\ tt}}$]}
    ]1[
    if $\mathcal{B}[\![b]\!]s = \mathbf{tt}$
    ]{
    \langle \mathtt{if}\ b\ \mathtt{then}\ S_1\ \mathtt{else}\ S_2, s \rangle \to s'
  }
\end{prooftree}
            \)
    \end{rightcolumn}
    \vspace{3em}
    \switchcolumn[1]*[]
    \begin{leftcolumn}
        \begin{Lean_Code}
  | if_false {b S@$_{1}$@ S@$_{2}$@ s s'}
    (h_cond : @\(\mathcal{B}[\![b]\!]\)@ s = false)
    (h_else : big_step S@$_{2}$@ s s') :
      big_step (Stmt.cond b S@$_{1}$@ S@$_{2}$@) s s'
        \end{Lean_Code}
    \end{leftcolumn}
    \begin{rightcolumn}
        \vspace{1.1em}
            \(
\begin{prooftree}
  \hypo{\langle S_2, s \rangle \to s'}
  \infer[
    left label={[if$_{\text{ns}}^{\text{\ ff}}$]}
    ]1[
    if $\mathcal{B}[\![b]\!]s = \mathbf{ff}$
    ]{
    \langle \mathtt{if}\ b\ \mathtt{then}\ S_1\ \mathtt{else}\ S_2, s \rangle \to s'
  }
\end{prooftree}
            \)
    \end{rightcolumn}
    \vspace{3em}
    \switchcolumn[1]*[]
    \begin{leftcolumn}
        \begin{Lean_Code}
  | while_true {b S s s' s''}
    (h_cond : @\(\mathcal{B}[\![b]\!]\)@ s = true)
    (h_step : big_step S s s')
    (h_rest : big_step (Stmt.loop b S) s' s'') :
      big_step (Stmt.loop b S) s s''
        \end{Lean_Code}
    \end{leftcolumn}
    \begin{rightcolumn}
        \vspace{1.1em}
            \(
\begin{prooftree}
  \hypo{\langle S, s \rangle \to s'}
  \hypo{\langle \mathtt{while}\ b\ \mathtt{do}\ S, s' \rangle \to s''}
  \infer[
      left label={[while$_{\text{ns}}^{\text{\ tt}}$]}
      ]2[
      if $\mathcal{B}[\![b]\!]s = \mathbf{tt}$
      ]{
    \langle \mathtt{while}\ b\ \mathtt{do}\ S, s \rangle \to s''
  }
\end{prooftree}
            \)
    \end{rightcolumn}
    \vspace{2em}
    \switchcolumn[1]*[]
    \begin{leftcolumn}
        \begin{Lean_Code}
  | while_false {b S' s}
    (h_cond : @\(\mathcal{B}[\![b]\!]\)@ s = false) :
      big_step (Stmt.loop b S') s s

-- Define the standard Nielson & Nielson notation
notation:40 "@\color{purple}$\langle$@" S "," s "@\color{purple}$\rangle$@" @\leanbreak@ "@\ \color{purple}\to$_{ns}$\ @" s' => @\leanbreak@ big_step S s s'
        \end{Lean_Code}
    \end{leftcolumn}
    \begin{rightcolumn}
        \vspace{1.1em}
            \(
\begin{prooftree}
  \hypo{} % Empty hypothesis for an axiom
  \infer[
    left label={[while$_{\text{ns}}^{\text{\ ff}}$]}
    ]1[
    if $\mathcal{B}[\![b]\!]s = \mathbf{ff}$
    ]{
    \langle \mathtt{while}\ b\ \mathtt{do}\ S, s \rangle \to s
  }
\end{prooftree}
            \)
    \end{rightcolumn}
\end{paracol}
\vspace{1em}


%
%
%
\section{Formalizing Proofs in Natural Semantics}
\label{formal_ns_sos_equiv:sec:ns:proofs}
Proofs made for natural semantics in the literature tend to share a pattern in using induction over the shape of the tree.
More formally the book uses the following definition:
\begin{quotation}
\noindent
Induction on the Shape of Derivation Trees\vspace{1em} \\  \noindent
1: Prove that the property holds for all the simple derivation trees by
showing that it holds for the axioms of the transition system.\vspace{1em} \\ \noindent
2: Prove that the property holds for all composite derivation trees: For
each rule assume that the property holds for its premises (this is
called the induction hypothesis) and prove that it also holds for the
conclusion of the rule provided that the conditions of the rule are
satisfied.
\end{quotation}
Doing induction over the shape will be used in the proof in \autoref{formal_ns_sos_equiv:sec:ns:while_unroll}
Let's start with an easier proof without any induction.

%
%
\subsection{While-loop Unrolling}
\label{formal_ns_sos_equiv:sec:ns:while_unroll}
In the textbook they provide a lemma stating the following:
\begin{quote}
The statement
\[\mathtt{while}\ b\ \mathtt{do}\ S\]
is semantically equivalent to
\[\mathtt{if}\ b\ \mathtt{then} (S;\ \mathtt{while}\ b\ \mathtt{do}\ S)\ \mathtt{else}\ \mathtt{skip}\]
\end{quote}
This would be equivalent with stating the following:
\[
    \langle \mathtt{while}\ b\ \mathtt{do}\ S,\ s\rangle \to s'' \iff  \langle \mathtt{if}\ b\ \mathtt{then} (S;\ \mathtt{while}\ b\ \mathtt{do}\ S)\ \mathtt{else}\ \mathtt{skip},\ s\rangle \to s''
\]
Making sure to label them the same way they do in the textbook we get the following:
\[
\langle \mathtt{while}\ b\ \mathtt{do}\ S,\ s\rangle \to s'' \tag{*}
\]
\[
\langle \mathtt{if}\ b\ \mathtt{then} (S;\ \mathtt{while}\ b\ \mathtt{do}\ S)\ \mathtt{else}\ \mathtt{skip},\ s\rangle \to s'' \tag{**}
\]
In the textbook they divide the proof into to two parts, first part is to prove that given (*) that (**) holds,
the second part proves the reverse, namely proving that given (**) that (*) holds. Which is the normal way to prove equivalence of two statements.

Let us now compare the textbook proof with the mechanized equivalent for the case where we show that (*) holds, implies that (**) holds.
\begin{paracol}{2}
\begin{leftcolumn}
    \centering \textbf{Mechanized (Lean 4)}
\end{leftcolumn}
\begin{rightcolumn}
    \centering \textbf{Pen-and-Paper (Textbook)} \\\noindent
\end{rightcolumn}

\switchcolumn[1]*[]
\begin{leftcolumn}
    \begin{Lean_Code}
theorem while_expand
    (b : Bexp)
    (S : Stmt)
    (s s' : State) :
  (@\(\langle\)@Stmt.loop b S, s@\(\rangle\)@ @\(\to_{ns}\)@ s') @\to@
  (@\(\langle\)@Stmt.cond
      b
      (Stmt.composition
          S
          (Stmt.loop b S))
      Stmt.skip, s@\(\rangle\)@ @\(\to_{ns}\)@ s') := by

  -- Assume (*) holds
  intro h_while

  -- Split into cases depending
  -- on if condition is true
  -- or not.
  cases h_while with
    \end{Lean_Code}
\end{leftcolumn}
\begin{rightcolumn} \noindent
Because (*) holds, we know that we have a derivation tree T for it. It can
have one of two forms depending on whether it has been constructed using the
rule [while$_{\text{ns}}^{\text{\ tt}}$] or the axiom [while$_{\text{ns}}^{\text{\ ff}}$].
\end{rightcolumn}
\switchcolumn[1]*[]
\begin{leftcolumn}
    \begin{Lean_Code}
  -- Invert while statement to
  -- access condition, and T@\color{green!40!black}$_1$@ T@\color{green!40!black}$_2$@
  | while_true
    h_cond_true T@$_1$@ T@$_2$@ =>
    \end{Lean_Code}
\end{leftcolumn}
\begin{rightcolumn} \noindent
In the first case, the derivation tree T has the form
\[
    \begin{prooftree}
        \hypo{T_1}
        \hypo{T_2}
        \infer[
            ]2{
            \langle \mathtt{while}\ b\ \mathtt{do}\ S,\ s\rangle \to s''
            }
    \end{prooftree}
\]
where \(T_1\) is a derivation tree with root \(\langle S,\ s\rangle \to s'\) and \(T_2\) is a derivation tree
with root \(\langle \mathtt{while}\ b\ \mathtt{do}\ S,\ s\rangle \to s''\).
Furthermore, \(\mathcal{B}[\![b]\!]s = \mathbf{tt}\).
\end{rightcolumn}
\vspace{1.5em}

\switchcolumn[1]*[]
\begin{leftcolumn}
    \begin{Lean_Code}
    -- Build the composition
    have comp_stmt :
      (@\(\langle\)@Stmt.composition S (Stmt.loop b S), s@\(\rangle\)@ @\(\to_{ns}\)@ s') :=
      -- Use comp rule with T@$_1$@ T@$_2$@ 
      big_step.comp T@$_1$@ T@$_2$@
    \end{Lean_Code}
\end{leftcolumn}

\begin{rightcolumn} \noindent
Using the derivation trees \(T_1\) and \(T_2\) as the premises for the rules [comp$_{\text{ns}}$], we can construct the derivation tree
\[
    \begin{prooftree}
        \hypo{T_1}
        \hypo{T_2}
        \infer[
            % left label=[comp$_{\text{ns}}$]
            ]2{
            \langle S;\ \mathtt{while}\ b\ \mathtt{do}\ S,\ s\rangle \to s''
            }
    \end{prooftree}
\]
\end{rightcolumn}
\vspace{1.5em}

\switchcolumn[1]*[]
\begin{leftcolumn}
    \begin{Lean_Code}
    -- Build if statement
    -- Use comp_stmt
    -- condition and if_true rule
    exact big_step.if_true
      h_cond_true -- Condition
      comp_stmt -- Then case
    \end{Lean_Code}
\end{leftcolumn}

\begin{rightcolumn} \noindent
Using that \(\mathcal{B}[\![b]\!]s = \mathbf{tt}\), we can use the rule [if$_{\text{ns}}^{\ \text{tt}}$]
to construct the derivation tree
\[
    \begin{prooftree}
        \hypo{T_1}
        \hypo{T_2}
        \infer[
            % left label=[comp$_{\text{ns}}$]
            ]2{
            \langle S;\ \mathtt{while}\ b\ \mathtt{do}\ S,\ s\rangle \to s''
            }
        \infer[
            % left label=[if$_{\text{ns}}^{\ \text{tt}}$]
            ]1{
            \langle \mathtt{if}\ b\ \mathtt{then}\ (S;\ \mathtt{while}\ b\ \mathtt{do}\ S)\ \mathtt{else}\ \mathtt{skip},\ s\rangle \to s''
            }
    \end{prooftree}
\]
hereby showing that (**) holds.
\end{rightcolumn}
\vspace{2em}
\switchcolumn[1]*[]
\begin{leftcolumn}
    \begin{Lean_Code}
  -- Invert while statement to
  -- access condition
  | while_false
    h_cond_false =>
    \end{Lean_Code}
\end{leftcolumn}
\begin{rightcolumn} \noindent
Alternatively, the derivation tree \(T\) is an instance of [while$_{\text{ns}}^{\ \text{ff}}$].
Then \(\mathcal{B}[\![b]\!]s = \mathbf{ff}\) and we must have that \(s'' = s\). So \(T\) simply is
\[
    \begin{prooftree}
        \hypo{\langle \mathtt{while}\ b\ \mathtt{do}\ S,\ s\rangle \to s}
    \end{prooftree}
\]
\end{rightcolumn}
\vspace{1.5em}

\switchcolumn[1]*[]
\begin{leftcolumn}
    \begin{Lean_Code}
    -- build skip statement
    have comp_stmt :
      (@\(\langle\)@Stmt.skip, s@\(\rangle\)@ @\(\to_{ns}\)@ s) :=
      -- Use axiom
      big_step.skip
    \end{Lean_Code}
\end{leftcolumn}
\begin{rightcolumn} \noindent
Using the axiom [skip$_{\text{ns}}$], we get a derivation tree
\[
    \begin{prooftree}
        \hypo{\langle \mathtt{skip},\ s\rangle \to s''}
    \end{prooftree}
\]
\end{rightcolumn}
\vspace{1.5em}

\switchcolumn[1]*[]
\begin{leftcolumn}
    \begin{Lean_Code}
    exact big_step.if_false
      h_cond_false -- Condition
      comp_stmt -- Else case
    \end{Lean_Code}
\end{leftcolumn}

\switchcolumn
\begin{rightcolumn} \noindent
and we can now apply the rule [if$_{\text{ns}}^{\ \text{ff}}$] to construct a derivation tree for (**):
\[
    \begin{prooftree}
        \hypo{\langle \mathtt{skip},\ s\rangle \to s''}
        \infer[
        ]1{
            \langle \mathtt{if}\ b\ \mathtt{then}\ (S;\ \mathtt{while}\ b\ \mathtt{do}\ S)\ \mathtt{else}\ \mathtt{skip},\ s\rangle \to s''
        }
    \end{prooftree}
\]
This completes the first part of the proof.
\end{rightcolumn}
\end{paracol}
\noindent\rule{\linewidth}{0.3pt}
As we can see the two version mimic each other very closely proving the same thing.
This proof is building the deduction tree from the top and moving down towards the root with our wanted result.
This method is also possible as demonstrated in Lean, but a more common approach is to build bottom up,
but it is entirely up to the writer which approach to use.

Now we will show that (**) holds, implies that (*) holds.
\begin{paracol}{2}
\begin{leftcolumn}
    \centering \textbf{Mechanized (Lean 4)}
\end{leftcolumn}
\begin{rightcolumn}
    \centering \textbf{Pen-and-Paper (Textbook)} \\\noindent
\end{rightcolumn}

\switchcolumn[1]*[]
\begin{leftcolumn}
    \begin{Lean_Code}
theorem if_compact (b : Bexp) (S : Stmt) (s s' : State) :
  (@\(\langle\)@Stmt.cond b (Stmt.composition S (Stmt.loop b S)) Stmt.skip, s@\(\rangle\)@ @\to$_{ns}$@ s') →
  (@\(\langle\)@Stmt.loop b S, s@\(\rangle\)@ @\to$_{ns}$@ s') := by
  
  -- Assume (**) holds
  intro h_if
    \end{Lean_Code}
\end{leftcolumn}
\begin{rightcolumn} \noindent
For the second part of the proof, we assume that (**) holds and shall prove
that (*) holds. So we have a derivation tree \(T\) for (**) and must construct one
for (*).
\end{rightcolumn}
\vspace{1.5em}

\switchcolumn[1]*[]
\begin{leftcolumn}
    \begin{Lean_Code}
  -- Split into cases depending
  -- on if condition is true
  -- or not.
  cases h_if with
    \end{Lean_Code}
\end{leftcolumn}
\begin{rightcolumn} \noindent
Only two rules could give rise to the derivation tree \(T\) for (**), namely
[while$_{\text{ns}}^{\text{\ tt}}$] or [while$_{\text{ns}}^{\text{\ ff}}$].
\end{rightcolumn}
\vspace{1.5em}

\switchcolumn[1]*[]
\begin{leftcolumn}
    \begin{Lean_Code}
  -- Invert if statement to
  -- access condition, and T@\color{green!40!black}$_1$@
  | if_true
    h_cond_true
    T@$_1$@ =>
    \end{Lean_Code}
\end{leftcolumn}
\begin{rightcolumn} \noindent
In the first case, \(\mathcal{B}[\![b]\!]s = \mathbf{tt}\) and we have a derivation tree \(T_1\)
with root
\[
    \begin{prooftree}
        \hypo{\langle S;\ \mathtt{while}\ b\ \mathtt{do}\ S,\ s\rangle \to s''}
    \end{prooftree}
\]
\end{rightcolumn}
\vspace{1.5em}

\switchcolumn[1]*[]
\begin{leftcolumn}
    \begin{Lean_Code}
    -- Deconstruct the composition derivation into its two parts.
    cases T@$_1$@ with
    | comp T@$_2$@ T@$_3$@ =>
    \end{Lean_Code}
\end{leftcolumn}
\begin{rightcolumn} \noindent
The statement has the general form \(S_1;\ S_2\), and the only rule that could give
this is [comp$_{\text{ns}}$]. Therefore there are derivation trees \(T _2\) and \(T_3\) for
\[
    \begin{prooftree}
        \hypo{\langle S,\ s\rangle \to s'}
    \end{prooftree}
\]
and
\[
    \begin{prooftree}
        \hypo{\langle \mathtt{while}\ b\ \mathtt{do}\ S,\ s'\rangle \to s''}
    \end{prooftree}
\]
for some state \(s'\).
\end{rightcolumn}
\vspace{1.5em}

\switchcolumn[1]*[]
\begin{leftcolumn}
    \begin{Lean_Code}
      exact big_step.while_true
        h_cond_true
        T@$_2$@ T@$_3$@
    \end{Lean_Code}
\end{leftcolumn}
\begin{rightcolumn} \noindent
It is now straightforward to use the rule [while$_{\text{ns}}^{\text{\ tt}}$] to combine
\(T _2\) and \(T_3\) into a derivation tree for (*).
\end{rightcolumn}
\vspace{1.5em}

\switchcolumn[1]*[]
\begin{leftcolumn}
    \begin{Lean_Code}
  | if_false
    h_cond_false
    h_deriv_skip =>
    \end{Lean_Code}
\end{leftcolumn}

\begin{rightcolumn} \noindent
In the second case, \(\mathcal{B}[\![b]\!]s = \mathbf{ff}\) and \(T\) is constructed using the rule [while$_{\text{ns}}^{\text{\ ff}}$].
\end{rightcolumn}
\vspace{1.5em}

\switchcolumn[1]*[]
\begin{leftcolumn}
    \begin{Lean_Code}
    -- The else branch is `skip`
    -- invert that derivation
    cases h_deriv_skip
    \end{Lean_Code}
\end{leftcolumn}

\begin{rightcolumn} \noindent
This means that we have a derivation tree for
\[
    \begin{prooftree}
        \hypo{\langle \mathtt{skip},\ s\rangle \to s''}
    \end{prooftree}
\]
and according to axiom [skip$_{\text{ns}}$] it must be the case that \(s = s''\).
\end{rightcolumn}
\vspace{1.5em}

\switchcolumn[1]*[]
\begin{leftcolumn}
    \begin{Lean_Code}
    -- produce `while_false`
    exact big_step.while_false
      h_cond_false
    \end{Lean_Code}
\end{leftcolumn}

\begin{rightcolumn} \noindent
But then we can use the axiom [while$_{\text{ns}}^{\text{\ ff}}$] to construct a derivation tree for (*).
This completes the proof.
\end{rightcolumn}
\vspace{1.5em}

\end{paracol} \noindent
\noindent\rule{\linewidth}{0.3pt}
Whilst the textbook is now done with thier proof, we must continue to formulate the equivalence relation,
this is done by using the previous two proofs in a new final proof constructued as following:
\begin{Lean_Code}
lemma while_unroll (b : Bexp) (S : Stmt) (s s' : State) :
(@\(\langle\)@(Stmt.cond
    b
    (Stmt.composition S (Stmt.loop b S))
    Stmt.skip), s@\(\rangle\)@ @\(\to_{ns}\)@ s') @\(\leftrightarrow\)@
(@\(\langle\)@Stmt.loop b S, s@\(\rangle\)@ @\(\to_{ns}\)@ s') := by
  -- equivalence introduction
  apply Iff.intro

  -- Forward direction: if => while
  case mp =>
    -- Reusing our previous proof for if => while
    exact if_compact b S s s'

  -- Backward direction: while => if
  case mpr =>
    -- Reusing our previous proof for while => if
    exact while_expand b S s s'
\end{Lean_Code}


%
%
\subsection{determinism}
\label{formal_ns_sos_equiv:sec:ns:determinism}
The next proof made in the textbook refers to determinism, where the definition of determinism is given
\begin{quotation}
\dots semantics is called deterministic if for all choices of
\(S\), \(s\), \(s'\) and \(s''\) we have that
\[
\langle S,\ s \rangle \to s'\ \mbox{and}\ \langle S,\ s \rangle \to s''\ \mbox{imply}\ s' = s''
\]
This means that for every statement \(S\) and initial state \(s\) we can uniquely
determine a final state \(s'\) if and only if the execution of \(S\) terminates.
\end{quotation}
By using this definition we can construct a proof using the following formulation:
\[
\langle S,\ s \rangle \to s' \implies \left(\langle S,\ s \rangle \to s'' \implies s' = s''\right)
\]
We will show the proofs side by side to show what the equivalent steps would be in Lean 4 compared to the textbook.
\vspace{1em}
\begin{paracol}{2}
    \begin{leftcolumn}
        \centering \textbf{Mechanized (Lean 4)}
    \end{leftcolumn}
    \begin{rightcolumn}
        \centering \textbf{Pen-and-Paper (Textbook)} \\\noindent
    \end{rightcolumn}
    \vspace{1.5em}


\switchcolumn[1]*[]
\begin{leftcolumn}
    \begin{Lean_Code}
theorem deterministic
  (S : Stmt)
  (s s' s'' : State)
  (h: ⟨S, s⟩ @\to$_{ns}$@ s') :
  (⟨ S, s⟩ @\to$_{ns}$@ s'')
  @\to@
  (s' = s'') := by
  -- Induction over the shape
  induction h
  -- by generalizing s'' we allow
  -- s'' to change across
  -- different induction steps
  generalizing s'' with
        \end{Lean_Code}
\end{leftcolumn}
\begin{rightcolumn}
We assume that \(\langle S,\ s \rangle \to s'\) and shall prove that
if \(\langle S,\ s \rangle \to s''\) then \(s' = s''\).
We shall proceed by induction on the shape of the derivation tree for \(\langle S,\ s \rangle \to s'\).
\end{rightcolumn}
\vspace{2em}


\switchcolumn[1]*[]
\begin{leftcolumn}
    \begin{Lean_Code}
  -- verifying tree built using
  -- [ass@\color{green!40!black}$_{ns}$@] axiom
  | ass =>
    -- Textbook:
    -- "only axiom or rule that
    -- could give ... is [ass@\color{green!40!black}$_{ns}$@]"

    -- Lean 4:
    -- Evaluate all matching
    -- cases by construction
    intro h_deriv_alt
    cases h_deriv_alt with
    | ass => rfl
    -- Lean only finds [ass@\color{green!40!black}$_{ns}$@] as
    -- the relevant case
    \end{Lean_Code}
\end{leftcolumn}
\begin{rightcolumn}
\textbf{The case} [ass$_{\text{ns}}$]: Then \(S\) is \(x := a\) and \(s'\) is \(s[x\mapsto\mathcal{A}[\![a]\!]s]\).
The only axiom or rule that could be used to give \( \langle x := a, s \rangle \to s'' \) is [ass$_{\text{ns}}$],
so it follows that \(s''\) must be \(s[x \mapsto \mathcal{A}[\![a]\!] s]\) and thereby \(s' = s''\).
\end{rightcolumn}
\vspace{2em}


\switchcolumn[1]*[]
\begin{leftcolumn}
    \begin{Lean_Code}
  -- same principle as above
  -- but with skip instead of ass
  | skip =>
    intro h_deriv_alt
    cases h_deriv_alt with
    | skip => rfl
    \end{Lean_Code}
\end{leftcolumn}
\begin{rightcolumn}
\textbf{The case} [skip$_{\text{ns}}$]: Analogous.
\end{rightcolumn}
\vspace{2em}


\switchcolumn[1]*[]
\begin{leftcolumn}
    \begin{Lean_Code}
  -- [comp@\color{green!40!black}$_{ns}$@] rule
  | comp
    h_left
    h_right
    ih_left
    ih_right =>
    intro h_deriv_alt
    cases h_deriv_alt with
    | comp
      h_alt_left
      h_alt_right =>
      -- textbook [compₙₛ]: first premises give equal intermediate state, then second premises give final equality
      have h_mid_eq := ih_left _ h_alt_left
      subst h_mid_eq
      exact ih_right _ h_alt_right
    \end{Lean_Code}
\end{leftcolumn}
\begin{rightcolumn}
\textbf{The case} [comp$_{\text{ns}}$]: Assume that
\[
\langle S_1;\ S_2,\ s \rangle \to s'
\]
holds because
\[
\langle S_1,\ s \rangle \to s_0\ \mbox{and}\ \langle S_2,\ s_0 \rangle \to s'
\]
for some \(s_0\). The only rule that could be applied to give \(\langle S_1;\ S_2,\ s \rangle \to s''\) is
[comp$_{\text{ns}}$], so there is a state \(s_1\) such that
\[
\langle S_1,\ s \rangle \to s_1\ \mbox{and}\ \langle S_2,\ s_1 \rangle \to s'
\]
The induction hypothesis can be applied to the premise \(\langle S_1,\ s \rangle \to s_0\) and from
\(\langle S_1,\ s \rangle \to s_1\) we get \(s_0 = s_1\). Similarly, the induction hypothesis can be applied
to the premise \(\langle S_2,\ s_0 \rangle \to s'\) and from \(\langle S_2,\ s \rangle \to s''\) we get \(s' = s''\) as required.
\end{rightcolumn}
\vspace{2em}


\switchcolumn[1]*[]
\begin{leftcolumn}
    \begin{Lean_Code}
  -- [if@\color{green!40!black}$_{\text{ns}}^{\ \text{tt}}$@] rule
  | if_true
    h_cond_true
    h_deriv_then
    ih_then =>
    intro h_deriv_alt
    cases h_deriv_alt with
    | if_true
      _ -- can ignore condition
      h_deriv_then_alt =>
      exact ih_then
        _
        h_deriv_then_alt

      -- Compared to the book
      -- we must consider the [if@\color{green!40!black}$_{\text{ns}}^{\ \text{ff}}$@] rule
      -- as we do not use side conditions in lean.
    | if_false
      h_cond_false_alt
      _ -- ignore statement
      =>
      -- We ignore the else stmt
      -- due to the trivial proof 
      -- from the contradiction
      -- of h_cond_true, and h_cond_false_alt.
      rw [h_cond_true] at h_cond_false_alt
      contradiction
    \end{Lean_Code}
\end{leftcolumn}
\begin{rightcolumn}
\textbf{The case} [if$_{\text{ns}}^{\ \text{tt}}$]: Assume that
\[
\langle \mathtt{if}\ b\ \mathtt{then}\ S_1\ \mathtt{else}\ S_2,\ s \rangle \to s'
\]
holds because
\[
\mathcal{B}[\![b]\!] = \mathbf{tt}\ \mbox{and}\ \langle S_1,\ s \rangle \to s'
\]
From \(\mathcal{B}[\![b]\!] = \mathbf{tt}\) we get that the only rule that could be applied to give the
alternative \(\langle \mathtt{if}\ b\ \mathtt{then}\ S_1\ \mathtt{else}\ S_2,\ s \rangle \to s''\) is 
[if$_{\text{ns}}^{\ \text{tt}}$]. So it must be the case that
\[
\langle S_1,\ s \rangle \to s'
\]
But then the induction hypothesis can be applied to the premise \(\langle S_1,\ s \rangle \to s'\)
and from \(\langle S_1,\ s \rangle \to s''\) we get \(s' = s''\).
\end{rightcolumn}
\vspace{2em}


\switchcolumn[1]*[]
\begin{leftcolumn}
    \begin{Lean_Code}
  -- same principle as above
  -- but with if_false instead of if_true
  -- and conditions being flipped
  -- [if@\color{green!40!black}$_{\text{ns}}^{\ \text{ff}}$@] rule
  | if_false
    h_cond_false
    h_deriv_else
    ih_else =>
    intro h_deriv_alt
    cases h_deriv_alt with
    | if_false
      _
      h_deriv_else_alt =>
      exact ih_else _ h_deriv_else_alt

    | if_true
      h_cond_true_alt
      _ =>
      rw [h_cond_false] at h_cond_true_alt
      contradiction
    \end{Lean_Code}
\end{leftcolumn}
\begin{rightcolumn}
\textbf{The case} [if$_{\text{ns}}^{\ \text{ff}}$]: Analogous.
\end{rightcolumn}
\vspace{2em}


\switchcolumn[1]*[]
\begin{leftcolumn}
    \begin{Lean_Code}
  -- [while@\color{green!40!black}$_{\text{ns}}^{\ \text{tt}}$@] rule
  | while_true
    h_cond_true
    h_deriv_body
    h_deriv_rest
    ih_while_body
    ih_while_rest =>
    intro h_deriv_alt
    cases h_deriv_alt with
    | while_false
      h_cond_false_alt =>
      rw [h_cond_true] at h_cond_false_alt
      contradiction
    | while_true
      _ -- can ignore condition
      h_deriv_body_alt
      h_deriv_rest_alt =>
      -- textbook [whileᵗᵗₙₛ]: align intermediate state via body IH, then apply loop IH
      have h_state_eq : _ = _ := ih_while_body _ h_deriv_body_alt
      subst h_state_eq
      exact ih_while_rest _ h_deriv_rest_alt
    \end{Lean_Code}
\end{leftcolumn}
\begin{rightcolumn}
\textbf{The case} [while$_{\text{ns}}^{\ \text{tt}}$]: Assume that
\[
\langle \mathtt{while}\ b\ \mathtt{do}\ S,\ s \rangle \to s'
\]
because
\[
\mathcal{B}[\![b]\!] = \mathbf{tt}, \ \langle S,\ s \rangle \to s_0\ \mbox{and}\ \langle \mathtt{while}\ b\ \mathtt{do}\ S,\ s_0 \rangle \to s'
\]
The only rule that could be applied to give \(\langle \mathtt{while}\ b\ \mathtt{do}\ S,\ s \rangle \to s''\)
is [while$_{\text{ns}}^{\ \text{tt}}$] because \(\mathcal{B}[\![b]\!] = \mathbf{tt}\), and this means that
\[
\langle S,\ s \rangle \to s_1\ \mbox{and}\ \langle \mathtt{while}\ b\ \mathtt{do}\ S,\ s_1 \rangle \to s''
\]
must hold for some \(s_1\). Again the induction hypothesis can be applied to the premise
\(\langle S,\ s \rangle \to s_0\), and from \(\langle S,\ s \rangle \to s_1\) we get \(s_0 = s_1\).
Thus we have
\[
\langle \mathtt{while}\ b\ \mathtt{do}\ S,\ s_0 \rangle \to s' \mbox{and}\ \langle \mathtt{while}\ b\ \mathtt{do}\ S,\ s_0 \rangle \to s''
\]
Since \(\langle \mathtt{while}\ b\ \mathtt{do}\ S,\ s_0 \rangle \to s'\) is a premise of [while$_{\text{ns}}^{\ \text{tt}}$],
we can apply the induction hypothesis to it. From \(\langle \mathtt{while}\ b\ \mathtt{do}\ S,\ s_0 \rangle \to s''\)
we therefore get \(s' = s''\) as required.
\end{rightcolumn}
\vspace{2em}


\switchcolumn[1]*[]
\begin{leftcolumn}
    \begin{Lean_Code}
  -- [while@\color{green!40!black}$_{\text{ns}}^{\ \text{ff}}$@] rule
  | while_false
    h_cond_false =>
    intro h_deriv_alt
    cases h_deriv_alt with

    -- Test creating from true rule
    | while_true
      h_cond_true_alt
      _ -- can ignore statment 
      _ -- can ignore remaining loops
      =>
      -- we have both that b is true, and false
      -- this is a contradiction
      rw [h_cond_false] at h_cond_true_alt
      contradiction

    | while_false
      _ -- we can ignore condition
      =>
      -- this is an axiom
      -- which is trivially true
      trivial
    \end{Lean_Code}
\end{leftcolumn}
\begin{rightcolumn}
\textbf{The case} [while$_{\text{ns}}^{\ \text{ff}}$]: Straightforward.
\end{rightcolumn}
\end{paracol}

%
%
%
\section{Formalizing Structural Operational Semantics (Small-Step)}
\label{formal_ns_sos_equiv:sec:sos:formalization}
Similar to the way we define natural semantics we can define structural operational semantics,
with the biggest diffrence being that not all statments will terminate to a final state after the first evaluation,but
instead might requier multiple steps to reach a final step for the operation. This is most notable for composiotion rules.

\vspace{1em}
\begin{paracol}{2}
    \begin{leftcolumn}
        \centering \textbf{Mechanized (Lean 4)}
    \end{leftcolumn}
    \begin{rightcolumn}
        \centering \textbf{Pen-and-Paper (Textbook)} \\\noindent
    \end{rightcolumn}

    \switchcolumn[1]*[]
    \begin{leftcolumn}
        \begin{Lean_Code}
inductive Config where
  | step : Stmt  @\to@ State  @\to@ Config
  | final : State  @\to@ Config
        \end{Lean_Code}
    \end{leftcolumn}
    \begin{rightcolumn}
    \end{rightcolumn}
    \vspace{2em}


    \switchcolumn[1]*[]
    \begin{leftcolumn}
        \begin{Lean_Code}
inductive small_step : Stmt @\to@ State @\to@ Config @\to@ Prop @\leanbreak@ where
        \end{Lean_Code}
    \end{leftcolumn}
    \begin{rightcolumn}
    \end{rightcolumn}
    \vspace{2em}


    \switchcolumn[1]*[]
    \begin{leftcolumn}
        \begin{Lean_Code}
  | ass {s x a} :
      small_step (Stmt.ass x a) s (Config.final (s[x @\mapsto@ @\(\mathcal{A}[\![a]\!]\)@ s]))
\end{Lean_Code}
    \end{leftcolumn}
    \begin{rightcolumn}
        \vspace{1.1em}
            \(
\begin{prooftree}
  \hypo{} % Empty hypothesis for an axiom
  \infer[left label=[ass$_{\text{sos}}$]]1{ \langle x := a, s \rangle \to s[x \mapsto \mathcal{A}[\![a]\!] s]}
\end{prooftree}
            \)
    \end{rightcolumn}
    \vspace{1em}


    \switchcolumn[1]*[]
    \begin{leftcolumn}
        \begin{Lean_Code}
  | skip {s} :
      small_step Stmt.skip s (Config.final s)
        \end{Lean_Code}
    \end{leftcolumn}
    \begin{rightcolumn}
        \vspace{1.1em}
            \(
\begin{prooftree}
  \hypo{} % Empty hypothesis for an axiom
  \infer[left label=[skip$_{\text{sos}}$]]1{ \langle \mathtt{skip}, s \rangle \to s}
\end{prooftree}
            \)
    \end{rightcolumn}
    \vspace{1em}
    \switchcolumn[1]*[]
    \begin{leftcolumn}
        \begin{Lean_Code}
  | comp1 {S@$_{1}$@ S@$_{1}$@' S@$_{2}$@ s s'}
    (h_step  : small_step S@$_{1}$@ s  (Config.step S@$_{1}$@' s')) :
      small_step (Stmt.composition S@$_{1}$@ S@$_{2}$@) s (Config.step (Stmt.composition S@$_{1}$@' S@$_{2}$@) s')
        \end{Lean_Code}
    \end{leftcolumn}
    \begin{rightcolumn}
        \vspace{1.1em}
            \(
\begin{prooftree}
  \hypo{\langle S_1, s \rangle \to \langle S_1', s' \rangle}
  \infer[
    left label=[comp$_{\text{sos}}^{1}$]
    ]1{
        \langle S_1;\ S_2, s \rangle \to \langle S_1';\ S_2, s' \rangle
    }
\end{prooftree}
            \)
    \end{rightcolumn}
    \vspace{3em}
    \switchcolumn[1]*[]
    \begin{leftcolumn}
        \begin{Lean_Code}
  | comp2 {S@$_{1}$@ S@$_{2}$@ s s'}
    (h_left  : small_step S@$_{1}$@ s (Config.final s')) :
      small_step (Stmt.composition S@$_{1}$@ S@$_{2}$@) s (Config.step S@$_{2}$@ s')
        \end{Lean_Code}
    \end{leftcolumn}
    \begin{rightcolumn}
        \vspace{1.1em}
            \(
\begin{prooftree}
  \hypo{\langle S_1, s \rangle \to s'}
  \infer[
    left label=[comp$_{\text{sos}}^{2}$]
  ]1{
    \langle S_1;\ S_2, s \rangle \to \langle S_2, s'\rangle}
\end{prooftree}
            \)
    \end{rightcolumn}
    \vspace{3em}
    \switchcolumn[1]*[]
    \begin{leftcolumn}
        \begin{Lean_Code}
  | if_true {b S@$_{1}$@ S@$_{2}$@ s s'}
    (h_cond_true : @\(\mathcal{B}[\![b]\!]\)@ s = true) : 
      small_step (Stmt.cond b S@$_{1}$@ S@$_{2}$@) s (Config.step S@$_{1}$@ s)
        \end{Lean_Code}
    \end{leftcolumn}
    \begin{rightcolumn}
        \vspace{1.1em}
            \(
\begin{prooftree}
  \hypo{} % Empty for axionm
  \infer[
    left label={[if$_{\text{sos}}^{\text{\ tt}}$]}
    ]1[
    if $\mathcal{B}[\![b]\!]s = \mathbf{tt}$
    ]{
        \parbox{4cm}{$
    \langle \mathtt{if}\ b\ \mathtt{then}\ S_1\ \mathtt{else}\ S_2, s \rangle \to \langle S_1, s\rangle
        $}
    }
\end{prooftree}
            \)
    \end{rightcolumn}
    \vspace{3em}
    \switchcolumn[1]*[]
    \begin{leftcolumn}
        \begin{Lean_Code}
  | if_false {b S@$_{1}$@ S@$_{2}$@ s s'}
    (h_cond_false : @\(\mathcal{B}[\![b]\!]\)@ s = false) : 
      small_step (Stmt.cond b S@$_{1}$@ S@$_{2}$@) s (Config.step S@$_{2}$@ s)
        \end{Lean_Code}
    \end{leftcolumn}
    \begin{rightcolumn}
        \vspace{1.1em}
            \(
\begin{prooftree}
  \hypo{} % Empty for axionom
  \infer[
    left label={[if$_{\text{ns}}^{\text{\ ff}}$]}
    ]1[
    if $\mathcal{B}[\![b]\!]s = \mathbf{ff}$
    ]{
        \parbox{4cm}{$
    \langle \mathtt{if}\ b\ \mathtt{then}\ S_1\ \mathtt{else}\ S_2, s \rangle \to \langle S_2, s\rangle
        $}
    }
\end{prooftree}
            \)
    \end{rightcolumn}
    \vspace{3em}
    \switchcolumn[1]*[]
    \begin{leftcolumn}
        \begin{Lean_Code}
  | while {b S s s' s''}
    (h_cond : @\(\mathcal{B}[\![b]\!]\)@ s = true)
    (h_step : big_step S s s')
    (h_rest : big_step (Stmt.loop b S) s' s'') :
      big_step (Stmt.loop b S) s s''

-- Define the standard Nielson & Nielson notation
-- Non terminal configuration
notation:40 "@\color{purple}$\langle$@" S "," s "@\color{purple}$\rangle$@" @\leanbreak@ "@\ \color{purple}\to$_{sos}$\ @" "@\color{purple}$\langle$@" S' "," s':40 "@\color{purple}$\rangle$@" => @\leanbreak@ small_step S s (Config.step S' s')

-- Terminal configuration
notation:40 "@\color{purple}$\langle$@" S "," s "@\color{purple}$\rangle$@" @\leanbreak@ "@\ \color{purple}\to$_{sos}$\ @" s':40 => @\leanbreak@ small_step S s (Config.final s')

        \end{Lean_Code}
    \end{leftcolumn}
    \begin{rightcolumn}
        \vspace{1.1em}
            \(
\begin{prooftree}
  \hypo{}
  \infer[
      left label={[while$_{\text{sos}}$]}
      ]1{
        \parbox{4cm}{$
            \langle \mathtt{while}\ b\ \mathtt{do}\ S, s \rangle \to \\
            \langle \mathtt{if}\ b\ \mathtt{then}\ (S;\ \mathtt{while}\ b\ \mathtt{do}\ S)\ \mathtt{else}\ \mathtt{skip}, s \rangle
        $}
    }
\end{prooftree}
            \)
    \end{rightcolumn}
\end{paracol}
\vspace{1em}

What we have defined is taking a single step using structural operationla semantics, but this can become cumbersome to reason with,
so the textbook added notation for taking some \(k\) steps using the notation \(\langle S, s\rangle \Rightarrow^{k} s'\),
and taking a finite amount of steps using the notation \(\langle S, s\rangle \Rightarrow^{*} s'\),
where the star is a unknown, but finite amount of steps,
this is used when the amount of steps is not important for the validity of the proof.
To use the same ideas in Lean 4, we need to introduce two new inductive types as follows:

\begin{Lean_Code}
inductive small_step_k : Config @\to@ Nat @\to@ Config @\to@ Prop where
  | refl {@\gamma@} :
      small_step_k @\gamma@ 0 @\gamma@
  | step {S s @\gamma@' @\gamma@'' k}
    (step : small_step S s @\gamma@')
    (rest : small_step_k @\gamma@' k @\gamma@'') :
      small_step_k (Config.step S s) (k + 1) @\gamma@''

-- Define the standard Nielson & Nielson notation
-- Non terminal configuration
notation:40 "@\color{purple}$\langle$@" S "," s "@\color{purple}$\rangle$@" @\leanbreak@ "@\ \color{purple}\to$_{sos}[$\ @" k "@\ \color{purple}$]$\ @" "@\color{purple}$\langle$@" S' "," s':40 "@\color{purple}$\rangle$@" => @\leanbreak@ small_step_k (Config.step S s) k (Config.step S' s')

-- Terminal configuration
notation:40 "@\color{purple}$\langle$@" S "," s "@\color{purple}$\rangle$@" @\leanbreak@ "@\ \color{purple}\to$_{sos}[$\ @" k "@\ \color{purple}$]$\ @" s':40 => @\leanbreak@ small_step_k (Config.step S s) k (Config.final s')

\end{Lean_Code}
and
\begin{Lean_Code}
inductive small_step_star : Config @\to@ Nat @\to@ Config @\to@ Prop where
  | refl {@\gamma@} :
      small_step_star @\gamma@ 0 @\gamma@
  | step {S s @\gamma@' @\gamma@'' k}
    (step : small_step S s @\gamma@')
    (rest : small_step_star @\gamma@' @\gamma@'') :
      small_step_star (Config.step S s) @\gamma@''

-- Define the standard Nielson & Nielson notation
-- Non terminal configuration
notation:40 "@\color{purple}$\langle$@" S "," s "@\color{purple}$\rangle$@" @\leanbreak@ "@\ \color{purple}\to$_{sos}*$\ @" "@\color{purple}$\langle$@" S' "," s':40 "@\color{purple}$\rangle$@" => @\leanbreak@ small_step_star (Config.step S s) (Config.step S' s')

-- Terminal configuration
notation:40 "@\color{purple}$\langle$@" S "," s "@\color{purple}$\rangle$@" @\leanbreak@ "@\ \color{purple}\to$_{sos}*$\ @" s':40 => @\leanbreak@ small_step_star (Config.step S s) (Config.final s')

\end{Lean_Code}

%
%
%
\section{Proving Natural Semantics $\equiv$ Structural Operational Semantics}
\label{formal_ns_sos_equiv:sec:ns_equiv_sos}
\todo{Add text}
\RED{\lipsum[1]}

%
\subsection{Formalization in Lean}
\label{formal_ns_sos_equiv:subsec:ns_equiv_sos:formalization_lean}
\todo{Add text}
\RED{\lipsum[1]}

%
\subsection{Challenges and Solutions}
\label{formal_ns_sos_equiv:subsec:ns_equiv_sos:chall_sol}
\todo{Add text}
\RED{\lipsum[1]}

%
\subsection{Insights Gained}
\label{formal_ns_sos_equiv:subsec:ns_equiv_sos:insights}
\todo{Add text}
\RED{\lipsum[1]}
